\newcommand{\tipoRequerimento}{Cobrança indevida sobre enquadramento tarifário incorreto}
\section{DOS FATOS E DOS DIREITOS}

\begin{enumerate}
    \item Em resposta a reclamação nº 003/2026 (Protocolo nº 1805222090), no qual se pleiteava o correto enquadramento da Unidade Consumidora (UC) nº 7133320 na subclasse B4a -- Iluminação Pública, a concessionária informou que a referida UC já se encontra devidamente classificada na referida tarifa.
    
    \item O reconhecimento técnico e administrativo por parte da distribuidora acerca do enquadramento em Iluminação Pública torna incontroversa a natureza da carga e a destinação da energia elétrica na unidade, qual seja, o atendimento estrito a áreas públicas de lazer e esporte de livre acesso.
    
    \item Contudo, a mera informação de que a unidade "já se encontra classificada" na subclasse IP omite fato primordial ao direito do consumidor: a data exata em que tal alteração cadastral foi efetivada e se houve a devida apuração e devolução dos valores faturados incorretamente no período pretérito.
    
    \item Conforme preconiza a regulação vigente da ANEEL, o reconhecimento do enquadramento incorreto atrai a obrigação imediata e compulsória da distribuidora de restituir os valores cobrados a maior durante todo o período em que a classificação inadequada persistiu, sob pena de enriquecimento ilícito da concessionária.
    
    \item Ademais, o direito à restituição retroativa observa o prazo prescricional de 10 (dez) anos, conforme pacificado pelo Despacho nº 2.006/2024 da ANEEL e fulcrado no art. 205 do Código Civil Brasileiro, devendo a devolução ocorrer em dobro, acrescida de atualização monetária pelo IPCA e juros de mora, haja vista a ausência de enquadramento tempestivo por parte da concessionária.
\end{enumerate}

\section{DOS PEDIDOS}

\begin{enumerate}
    \item Diante do exposto, e considerando o reconhecimento administrativo do enquadramento da UC na subclasse IP, requer que a Concessionária:
    
    \begin{enumerate}
        \item Apresente formalmente o histórico cadastral da UC nº 7133320, indicando de forma clara a data exata em que foi realizada a alteração para a tarifa B4a -- Iluminação Pública.
        
        \item Informe de maneira expressa se já foi realizada, de ofício, a compensação financeira pelos períodos em que a tarifação errada foi aplicada na unidade consumidora.
        
        \item Caso a compensação já tenha sido realizada, requer o envio imediato do respectivo memorial de cálculo detalhado do passivo, bem como do comprovante de crédito ou ordem de pagamento emitida em favor do Município.
        
        \item Caso a compensação retroativa não tenha sido efetuada, requer que a distribuidora proceda com a imediata apuração e devolução de todos os valores cobrados a maior relativos aos últimos 10 (dez) anos antecedentes ao correto enquadramento, em observância ao art. 205 do Código Civil e ao Despacho nº 2.006 da ANEEL.
        
        \item Que a referida devolução dos valores retroativos seja feita em dobro, atualizada monetariamente pelo IPCA e acrescida de juros de mora de 1\% ao mês calculados \textit{pro rata die}, conforme determina a Resolução Normativa nº 1.000/2021 da ANEEL.
        
        \item Proceda com a devolução do valor integral apurado por meio de depósito bancário na conta corrente do Município, nos termos do § 7º do art. 323 da REN nº 1.000/2021 da ANEEL, comunicando por escrito o resultado desta reclamação por meio de e-mail institucional.
    \end{enumerate}
\end{enumerate}